\section{The cohomology used for the host and for patching}

LD2--3 and DX1--3 established the finite dualities and Euler
formulas used below. We match them to the host and patching
records, and compute the additional tame symbol in the cup order
required by the latter.

% BEGIN G71H1
\paragraph{G71H1. The coefficient dual and the local Euler factor.}
Let $F$ be a number field with the stipulated algebraic closure
$\Omega$. Fix $m=\ell^r$, $r>0$, $R=\mathbb Z/m$, and any finite
discrete $R$-module $M$ with the continuous action specified in
the source. A prolongation of a finite place $v$ identifies its
decomposition group with the absolute Galois group of $F_v$:
one first takes the henselization for this valuation and then its
completion, which has the same finite separable extensions. These
identifications carry the given inertia, residual Frobenius and
cyclotomic character to the corresponding local ones. The source
cohomology uses continuous inhomogeneous cochains. For any finite
$R$-module the map $x\mapsto x/m\bmod\mathbb Z$ on $R$ identifies
its $R$-linear dual with its Pontryagin dual. With a primitive
$m$th root $\zeta$, $x\mapsto\zeta^x$ identifies the cyclotomic
coefficient with $\mu_m$. Thus the source's twisted dual is
$M^D=\operatorname{Hom}(M,\mu_m)$ with its usual Galois action.

Use Tate local duality and the local Euler characteristic theorem,
in the forms proved in Milne, \emph{Arithmetic Duality Theorems},
second edition, Ch.~I, \S2. These give finiteness in degrees
zero through two, vanishing above two, and perfect cup pairings
$H^i(M)\times H^{2-i}(M^D)\to H^2(\mu_m)$. The local Brauer
invariant identifies the last group with $R$. Composing with
this isomorphism recovers the original bijectivity assertions for
the $R$-linear cup maps, in both orders for degree one. Write
$h_v^i(M)=|H^i(F_v,M)|$. The local Euler formula gives
\[
 h_v^1(M)=h_v^0(M)h_v^2(M)
 \begin{cases} |M|^{[F_v:\mathbb Q_\ell]} &v\mid\ell,\\
 1 &v\nmid\ell .\end{cases}
\]
At $v\mid\ell$, this degree is the ramification index times the
residue degree, the exact exponent $\mathrm{locDeg}(v)$.
These comparisons give the three fields of local finiteness,
Tate duality and Euler characteristic in both source packages.
% END G71H1

% BEGIN G71H2
\paragraph{G71H2. Global duality with the source's chosen invariant maps.}
For the finite set $S$ containing all places over $\ell$, the
source quotient $G_S$ kills the closed normal subgroup generated
by the inertias outside $S$. It is the Galois group of the maximal
extension unramified away from $S$ and infinity. Hence the finite
module in the question is unramified at every omitted finite
place. The global finiteness and Euler theorems give finite
$H^0,H^1,H^2$, and, for odd $\ell$, vanishing in higher degrees.
Writing $h^i=|H^i(G_S,M)|$, they give
\[
  h^0 h^2 |M|^{[F:\mathbb Q]}=h^1 A_\infty(M),\qquad
  A_\infty(M)=|M|^{r_2}\prod_{v\ {\rm real}}|M^{c_v}|,
\]
where $r_2$ counts complex places and $c_v$ is the specified
complex conjugation. The truncated Euler identity also holds at
$\ell=2$ for real fields. There the source imposes no general
vanishing in degrees above two. See Milne, Ch.~I,
Thms.~4.10 and 5.1.

Fix the primitive root $\zeta$ before choosing any $S$, module or
prolongations in the global test. At every valuation above a
finite place identify the cyclotomic coefficient with $\mu_m$
using this same $\zeta$. Identify the local Brauer $m$-torsion
with $R$ by the invariant, using the common sign in which
\[
  \operatorname{inv}_v(\kappa_\varpi\smile\chi_{\rm ur})=1,
  \qquad \chi_{\rm ur}(\operatorname{Frob}_{\rm ar})=1 .
\]
Indeed the Kummer class of a uniformizer paired with the
unramified degree-$m$ character gives the cyclic algebra of that
unramified extension and the uniformizer, up to the universal
cup-order sign. Its invariant has numerator one since
$v(\varpi)=1$. This works also at $v\mid\ell$. It is independent
of the chosen uniformizer and of its $m$th root. A homomorphism
trivial on inertia and taking value one on arithmetic Frobenius
is the specified unramified character: on the procyclic quotient
it factors through reduction modulo $m$. The source cocycle
$z(g,h)=\kappa_\varpi(g)\chi_{\rm ur}(h)$ therefore has the
displayed invariant. At valuations lying above no finite place
set the map to zero. This gives one family with the literal
$\mathrm{IsNormalised}$ test wherever required.

Under $\ell\ne2$ or $F$ totally imaginary, apply Poitou--Tate
duality for $S$ with the infinite places adjoined (Milne,
Ch.~I, Thm.~4.10). All modified infinite local groups then
vanish. Its nine terms are the global cohomology in degrees
zero to two, the products of the finite local groups, and the
duals of the complementary global groups for $M^D$. The three
outgoing local maps pair with the respective restrictions and
sum the local invariants, precisely $\mathrm{tr0},\mathrm{tr1},
\mathrm{tr2}$. Reversing degree-one cup order or changing the
common sign multiplies the corresponding map by $-1$ and
preserves both its kernel and its image. The connecting maps
give the two existential linear arrows in $\mathrm{NineTermExact}$.
The duality of the localization kernels gives, in either order,
$\Sha^2(M)\simeq\Sha^1(M^D)^*$, with both sides finite of the
same order. The higher global groups vanish under this same
archimedean guard. This proves $\mathrm{PoitouTateNormalised}$,
including the quantifier on the invariant family before all
choices of $S$ and $M$.
% END G71H2

% BEGIN G71H3
\paragraph{G71H3. The unramified orthogonal and the Selmer order identity.}
At $v\nmid\ell$ with $M$ unramified, $M^D$ is unramified too.
An unramified $H^1$ class inflates from the procyclic residue
Galois group. Two such classes have zero cup product since
that group has $\ell$-cohomological dimension one.
Moreover $H^1_{\rm ur}(M)=M/(\mathrm{Frob}-1)M$ has order
$h_v^0(M)$. The dual unramified group has order
$h_v^0(M^D)=h_v^2(M)$. By G71H1 their product is $h_v^1(M)$.
Perfectness gives the full mutual orthogonals.

For odd $\ell$ put $V=\prod_{v\in S}H^1(F_v,M)$ and let
$X\subset V$ be the image of global localization. By G71H2 the
image for $M^D$ is $X^\perp$ in the perfect dual product. For
arbitrary local conditions $D=\prod_{v\in S}D_v$ the source's
dual conditions are $D^\perp$. Taking the two inverse images
under localization gives
\[
 \frac{|\mathrm{Sel}_D(M)|}{|\mathrm{Sel}_{D^\perp}(M^D)|}
 =
 \frac{|\Sha^1(M)|}{|\Sha^1(M^D)|}
 \frac{|X\cap D|}{|(X+D)^\perp|}
 =
 \frac{h^1(M)\,|D|}{|\Sha^2(M)|\,|V|}.
\]
The last equality uses $|X+D||X\cap D|=|X||D|$ and kernel
duality. The last three terms of Poitou--Tate give
$|\Sha^2(M)|=h^2(M)h^0(M^D)/\prod_{v\in S}h_v^2(M)$.
Insert the local and global Euler identities. Since $S$ contains
all places above $\ell$, their local degrees sum to
$[F:\mathbb Q]$. The result is
\[
 \frac{|\mathrm{Sel}_D(M)|}{|\mathrm{Sel}_{D^\perp}(M^D)|}
 =\frac{h^0(M)\prod_{v\in S}|D_v|}
 {h^0(M^D)A_\infty(M)\prod_{v\in S}h_v^0(M)}.
\]
This is the original $\mathrm{GreenbergWiles}$ for any chosen
prolongations and every local condition. Together with G71H1--2
and unramified orthogonality it supplies all eight fields of
$\mathrm{GC.Hypotheses}_2$. Thus AFE16 now constructs the
specified marked host and its seed clause from the original host
data, using these classical duality theorems.
% END G71H3

% BEGIN G71H4
\paragraph{G71H4. The tame quotient and the signed local symbol.}
At $v\nmid\ell$ write $p$ for its residue characteristic and
$q$ for the residue cardinality. Wild inertia is pro-$p$ and
tame inertia is $\prod_{s\ne p}\mathbb Z_s(1)$. Therefore the
kernel $J$ of the maximal pro-$\ell$ quotient of inertia $I$ is
the inverse image of the factors $s\ne p,\ell$. It is closed,
normal in the decomposition group and all its finite quotients
have order prime to $\ell$. The quotient $I/J$ is procyclic
isomorphic to $\mathbb Z_\ell(1)$, with a generator lifting to
some $\gamma\in I$. An arbitrary arithmetic Frobenius lift
$\phi$ acts on $I/J$ as the $q$th power. The powers of $\phi$
extend to a section of the procyclic unramified quotient
$\widehat{\mathbb Z}$, giving the semidirect product description
of the decomposition group modulo $J$. Thus $\phi,\gamma$
generate modulo $J$ topologically. If $a,c$ in a finite group
satisfy $c^{\ell^k}=1$ and $aca^{-1}=c^q$, the two continuous
maps from the procyclic factors sending their generators to
$a,c$ give the required homomorphism with open kernel. This
proves every assertion in $\mathrm{TameQuotient}$ for any of
the specified lifts and generators.

Suppose $q\equiv1\bmod m$, $m=\ell^n$. The case $m=1$ has
the zero coefficient ring and all identities hold there.
Otherwise Hensel lifting gives $\mu_m\subset F_v$, so all
degree-one cocycles below are homomorphisms. For the chosen
primitive root let $\kappa$ be the Kummer character of the
specified $m$th root of a uniformizer, and let $\chi$ be the
unramified character with $\chi(\phi)=1$. The extension made
by this root is totally ramified of degree $m$ by Eisenstein,
so the restriction of $\kappa$ to $I/J$ is surjective. In
particular $\kappa$ takes a unit value on every pro-generator.
Use the local invariant in the order prescribed by the
separate $\mathrm{TameSymbol}$ test,
$\operatorname{inv}(\chi\smile\kappa)=1$. This normalization
exists by the cyclic-algebra calculation in G71H2. It gives
the required formula for every unramified character in the
first slot, for every uniformizer root and every Frobenius lift.

For any homomorphism $z$ to $R$, write $z=b\kappa+d\chi$
since $\kappa$ generates the tame direction. For the second
cocycle $y$, the unramified-class assumption means $y=y(\phi)\chi$,
as coboundaries vanish for these trivial coefficients.
The unramified cup square vanishes, and graded commutativity
gives $\operatorname{inv}(z\smile y)=-b\,y(\phi)$ for every
map satisfying the source normalization. At an arbitrary
inertial $\gamma$ the exponent $c$ determined by its action on
the root equals $\kappa(\gamma)$, so $z(\gamma)=b c$. Hence
\[
       c\,\operatorname{inv}(z\smile y)=-z(\gamma)y(\phi)
\]
also when $c$ is a nonunit. This is the requested sign and
order of the arguments. With G71H1--3 it fills the eighth
patching field: the full $\mathrm{Pa.CohP}$ holds on the
classical foundations used here. The tame presentation used
above is the standard one from the ramification theory of
local fields (Serre, \emph{Corps locaux}, Ch.~IV).
% END G71H4
